% (c) Nikita Lisitsa, lisyarus@gmail.com, 2026

\documentclass[10pt]{beamer}

\usepackage[T2A]{fontenc}
\usepackage[russian]{babel}
\usepackage{minted}

\usepackage[most]{tcolorbox}
\tcbuselibrary{minted}

\usepackage{graphicx}
\graphicspath{ {./images/} }

\usepackage{adjustbox}

\usepackage{color}
\usepackage{soul}
\usepackage{multicol}
\usepackage{hyperref}

\usetheme{metropolis}

\definecolor{red}{rgb}{1,0,0}
\definecolor{codebg}{RGB}{29,35,49}
\setminted{fontsize=\footnotesize}

\makeatletter
\newcommand{\slideimage}[1]{
  \begin{figure}
    \begin{adjustbox}{width=\textwidth, totalheight=\textheight-2\baselineskip-2\baselineskip,keepaspectratio}
      \includegraphics{#1}
    \end{adjustbox}
  \end{figure}
}
\makeatother

\title{Язык программирования C++}
\subtitle{Лекция 13: наследование, полиморфизм, стандартная библиотека, шаблоны}
\date{2026}
\author{Никита Лисица}

\setbeamertemplate{footline}[frame number]

\begin{document}

\frame{\titlepage}

% План:
% + Множественное наследование
% + Виртуальное наследование
% + Интерфейсы
% + Вложенные типы и классы
% + Неймспейсы
% + Using (type alias)
% + Стандартная библиотека
% * Пример: сортировка в стиле С
% * Пример: сортировка в ООП-стиле
% * Пример: связный список в стиле С
% * Пример: связный список в ООП-стиле
% * Пример: Array, scoped_ptr, ...
% * Шаблоны

\begin{frame}[fragile]
\frametitle{Множественное наследование}
\begin{itemize}
\usemintedstyle{lightbulb}
\begin{tcblisting}{listing engine=minted, minted language=cpp, minted options={fontsize=\scriptsize}, listing only, colback=codebg, colframe=codebg, rounded corners}
class Readable {
public:
  virtual void read(char* dst, size_t bytes) = 0;
};

class Writable {
public:
  virtual void write(const char* src, size_t bytes) = 0;
};
\end{tcblisting}
\pause
\usemintedstyle{tango}
\item Хотим сделать класс, который является и \mintinline{cpp}|Readable|, и \mintinline{cpp}|Writable|, но как?
\end{itemize}
\end{frame}

\begin{frame}[fragile]
\frametitle{Множественное наследование}
\begin{itemize}
\item С помощью \textit{множественного наследования (multiple inheritance)}:
\usemintedstyle{lightbulb}
\begin{tcblisting}{listing engine=minted, minted language=cpp, minted options={fontsize=\scriptsize}, listing only, colback=codebg, colframe=codebg, rounded corners}
class File : public Readable, public Writable {
public:
  void read(char* dst, size_t bytes) override { ... }
  void write(const char* src, size_t bytes) override { ... }
};
\end{tcblisting}
\pause
\usemintedstyle{tango}
\item Такой класс содержит все поля и методы всех классов, от которых наследуется
\pause
\item Указатель (или ссылка) такой класс приводится к типу указателя (или ссылки) на любой из базовых классов
\end{itemize}
\end{frame}

\begin{frame}[fragile]
\frametitle{Множественное наследование}
\begin{itemize}
\item Что, если у базовых классов есть одинаковые поля/методы?
\usemintedstyle{lightbulb}
\begin{multicols}{2}
\begin{tcblisting}{listing engine=minted, minted language=cpp, minted options={fontsize=\scriptsize}, listing only, colback=codebg, colframe=codebg, rounded corners}
class Readable {
protected:
  size_t totalBytes;
public:
  ...
};
\end{tcblisting}
\begin{tcblisting}{listing engine=minted, minted language=cpp, minted options={fontsize=\scriptsize}, listing only, colback=codebg, colframe=codebg, rounded corners}
class Writable {
protected:
  size_t totalBytes;
public:
  ...
};
\end{tcblisting}
\end{multicols}
\pause
\usemintedstyle{tango}
\item Можно использовать явное разрешение имён:
\usemintedstyle{lightbulb}
\begin{tcblisting}{listing engine=minted, minted language=cpp, minted options={fontsize=\scriptsize}, listing only, colback=codebg, colframe=codebg, rounded corners}
class File : public Readable, public Writable {
public:
  size_t bytesRead() const { return Readable::totalBytes; }
  size_t bytesWritten() const { return Writable::totalBytes; }
};
\end{tcblisting}
\end{itemize}
\end{frame}

\begin{frame}[fragile]
\frametitle{Виртуальное наследование}
\begin{itemize}
\item Что будет, если оба базовых класса сами наследуются от одного и того же класса (т.н. \textit{diamond problem})?
\usemintedstyle{lightbulb}
\begin{tcblisting}{listing engine=minted, minted language=cpp, minted options={fontsize=\scriptsize}, listing only, colback=codebg, colframe=codebg, rounded corners}
class IOBase {
protected:
  size_t totalBytes;
};

class Readable : public IOBase {
  ...
};

class Writable : public IOBase {
  ...
};

class File : public Readable, public Writable {
  ...
};
\end{tcblisting}
\end{itemize}
\end{frame}

\begin{frame}[fragile]
\frametitle{Виртуальное наследование}
\begin{itemize}
\usemintedstyle{tango}
\item У обоих базовых классов будет своя копия подобъекта \mintinline{cpp}|IOBase|
\pause
\item \(\Longrightarrow\) В классе \mintinline{cpp}|File| содержатся два подобъекта типа \mintinline{cpp}|IOBase|
\pause
\item Иногда это ровно то, что нужно; а иногда хочется, чтобы подобъект общего типа был один
\pause
\item Решение: \textit{виртуальное наследование (virtual inheritance)}:
\usemintedstyle{lightbulb}
\begin{tcblisting}{listing engine=minted, minted language=cpp, minted options={fontsize=\scriptsize}, listing only, colback=codebg, colframe=codebg, rounded corners}
class Readable : public virtual IOBase {
  ...
};

class Writable : public virtual IOBase {
  ...
};
\end{tcblisting}
\end{itemize}
\end{frame}

\begin{frame}[fragile]
\frametitle{Виртуальное наследование}
\begin{itemize}
\usemintedstyle{tango}
\item При виртуальном наследовании подобъект типа \mintinline{cpp}|IOBase| будет только один
\pause
\item Почему это свойство классов \mintinline{cpp}|Readable| и \mintinline{cpp}|Writable|, а не класса \mintinline{cpp}|File|?
\pause
\item \mintinline{cpp}|Readable/Writable| не могут независимо хранить подобъект \mintinline{cpp}|IOBase|, и вместо этого хранят \textit{указатель} на подобъект \mintinline{cpp}|IOBase|
\pause
\item Сам подобъект \mintinline{cpp}|IOBase| кладётся где-то по усмотрению компилятора, часто -- в конце класса
\end{itemize}
\end{frame}

\begin{frame}[fragile]
\frametitle{Виртуальное наследование}
\begin{itemize}
\usemintedstyle{tango}
\item Расположение данных в памяти получится аналогичным вот такому коду:
\usemintedstyle{lightbulb}
\begin{tcblisting}{listing engine=minted, minted language=cpp, minted options={fontsize=\scriptsize}, listing only, colback=codebg, colframe=codebg, rounded corners}
class IOBase { ... };

class Readable {
  // Аналог public virtual IOBase
  IOBase* ioBasePtr;
};

class Writable {
  // Аналог public virtual IOBase
  IOBase* ioBasePtr;
};

class File : public Readable, public Writable {
  // Компилятор сам положит подобъект IOBase
  // и правильно запишет указатели ioBasePtr
  IOBase ioBase;
}
\end{tcblisting}
\end{itemize}
\end{frame}

\begin{frame}[fragile]
\frametitle{Виртуальное наследование}
\begin{itemize}
\usemintedstyle{tango}
\item Конструктор виртуального базового класса должен явно вызывать любой класс, наследующийся от него (напрямую или через другие классы):
\usemintedstyle{lightbulb}
\begin{tcblisting}{listing engine=minted, minted language=cpp, minted options={fontsize=\scriptsize}, listing only, colback=codebg, colframe=codebg, rounded corners}
class IOBase {
public:
  IOBase(int flags);
};

class Readable : public virtual IOBase {
public:
  Readable() : IOBase(READ_FLAG) {}
};

class Writable : public virtual IOBase {
  Writable() : IOBase(WRITE_FLAG) {}
};

class File : public Readable, public Writable {
public:
  File() : IOBase(READ_FLAG | WRITE_FLAG) {}
};
\end{tcblisting}
\end{itemize}
\end{frame}

\begin{frame}[fragile]
\frametitle{Интерфейсы}
\begin{itemize}
\usemintedstyle{tango}
\item ООП-языки часто отдельно выделяют \textit{интерфейсы}: спецификации наборов методов, которые должен поддерживать некоторый класс
\pause
\item В C++ роль интерфейсов выполняют классы с \textit{чисто виртуальными методами} и \textit{без полей}:
\usemintedstyle{lightbulb}
\begin{tcblisting}{listing engine=minted, minted language=cpp, minted options={fontsize=\scriptsize}, listing only, colback=codebg, colframe=codebg, rounded corners}
class Hashable {
public:
  virtual size_t getHash() = 0;

  // Не забываем про виртуальный деструктор!
  virtual ~Hashable() {}
};

class Comparable {
public:
  virtual bool equals(Comparable* other) = 0;

  // Не забываем про виртуальный деструктор!
  virtual ~Comparable() {}
};
\end{tcblisting}
\end{itemize}
\end{frame}

\begin{frame}[fragile]
\frametitle{Интерфейсы}
\begin{itemize}
\usemintedstyle{tango}
\item Обычно лучше использовать множественное наследование только от интерфейсов -- некоторые ООП-языки (Java, C\#) разрешают только такое множественное наследование
\pause
\item Виртуальное наследование -- подход C++ к решению diamond problem, в других языках есть другие решения, см. \href{https://en.wikipedia.org/wiki/Multiple\_inheritance\#Mitigation}{\nolinkurl{en.wikipedia.org/wiki/Multiple\_inheritance\#Mitigation}}
\end{itemize}
\end{frame}

\begin{frame}[fragile]
\frametitle{Конфликты имён}
\begin{itemize}
\usemintedstyle{tango}
\item Часто возникает ситуация, когда в разных частях программы появились функции/классы с одинаковыми именами:
\begin{multicols}{2}
\usemintedstyle{lightbulb}
\begin{tcblisting}{listing engine=minted, minted language=cpp, minted options={fontsize=\scriptsize}, listing only, colback=codebg, colframe=codebg, rounded corners}
struct Node {
  int data;
  Node* next;
};


class List {
  Node* head;
  ...
};
\end{tcblisting}
\begin{tcblisting}{listing engine=minted, minted language=cpp, minted options={fontsize=\scriptsize}, listing only, colback=codebg, colframe=codebg, rounded corners}
struct Node {
  int data;
  Node* left;
  Node* right;
};

class Tree {
  Node* root;
  ...
};
\end{tcblisting}
\end{multicols}
\pause
\item Это приведёт к ошибкам компиляции, если какой-то код использует и \mintinline{cpp}|List|, и \mintinline{cpp}|Tree|
\pause
\item (Если неповезёт, может обойтись без ошибок компиляции, но с runtime ошибками!)
\end{itemize}
\end{frame}

\begin{frame}[fragile]
\frametitle{Конфликты имён}
\begin{itemize}
\usemintedstyle{tango}
\item Как избежать конфликтов имён? Вариант 1: префиксы в именах
\begin{multicols}{2}
\usemintedstyle{lightbulb}
\begin{tcblisting}{listing engine=minted, minted language=cpp, minted options={fontsize=\scriptsize}, listing only, colback=codebg, colframe=codebg, rounded corners}
struct ListNode {
  int data;
  ListNode* next;
};

\end{tcblisting}
\begin{tcblisting}{listing engine=minted, minted language=cpp, minted options={fontsize=\scriptsize}, listing only, colback=codebg, colframe=codebg, rounded corners}
struct TreeNode {
  int data;
  TreeNode* left;
  TreeNode* right;
};
\end{tcblisting}
\end{multicols}
\pause
\usemintedstyle{tango}
\item Длиннее имена функций/структур/классов, сложнее читать
\item Переводит проблему на уровень выше: кто-то мог создать два разных класса \mintinline{cpp}|Tree| :)
\end{itemize}
\end{frame}

\begin{frame}[fragile]
\frametitle{Конфликты имён}
\begin{itemize}
\usemintedstyle{tango}
\item Как избежать конфликтов имён? Вариант 2: вложенные типы
\begin{multicols}{2}
\usemintedstyle{lightbulb}
\begin{tcblisting}{listing engine=minted, minted language=cpp, minted options={fontsize=\scriptsize}, listing only, colback=codebg, colframe=codebg, rounded corners}
class List {
  struct Node {
    int data;
    Node* next;
  };

  Node* head;
};

\end{tcblisting}
\begin{tcblisting}{listing engine=minted, minted language=cpp, minted options={fontsize=\scriptsize}, listing only, colback=codebg, colframe=codebg, rounded corners}
class Tree {
  struct Node {
    int data;
    Node* left;
    Node* right;
  };

  Node* root;
};
\end{tcblisting}
\end{multicols}
\pause
\usemintedstyle{tango}
\item Отлично подходит в этом конкретном примере
\item Инкапсуляция: делает и сам тип \mintinline{cpp}|Node| недоступным снаружи
\item Не подходит, если нет класса, который логически должен содержать наши функции/классы
\end{itemize}
\end{frame}

\begin{frame}[fragile]
\frametitle{Пространства имён}
\begin{itemize}
\usemintedstyle{tango}
\item Как избежать конфликтов имён? Вариант 3: \textit{пространства имён (namespaces)}
\usemintedstyle{lightbulb}
\begin{tcblisting}{listing engine=minted, minted language=cpp, minted options={fontsize=\scriptsize}, listing only, colback=codebg, colframe=codebg, rounded corners}
namespace xml {
  class Node { ... };
  class Parser { ... };
}

namespace json {
  class Node { ... };
  class Parser { ... };
}

json::Parser parser("myFile.json");
...
\end{tcblisting}
\pause
\item При компиляции пространства имён тоже используются для name mangling'а
\end{itemize}
\end{frame}

\begin{frame}[fragile]
\frametitle{Пространства имён}
\begin{itemize}
\usemintedstyle{tango}
\item Если пространство имён используется часто, \textit{можно} подключить его с помощью конструкции \mintinline{cpp}|using namespace|:
\usemintedstyle{lightbulb}
\begin{tcblisting}{listing engine=minted, minted language=cpp, minted options={fontsize=\scriptsize}, listing only, colback=codebg, colframe=codebg, rounded corners}
using namespace xml;
Parser parser("myFile.xml");
\end{tcblisting}
\pause
\usemintedstyle{tango}
\item Лучше так не делать без крайней необходимости: это усложняет чтение кода и может приводить к конфликтам имён
\pause
\item В заголовочных файлах точно не стоит делать \mintinline{cpp}|using namespace|: оно неявно добавится во все файлы, использующие этот заголовочный файл
\pause
\item Совет: \textbf{не делайте} \mintinline{cpp}|using namespace std;|
\end{itemize}
\end{frame}

\begin{frame}[fragile]
\frametitle{Пространства имён}
\begin{itemize}
\usemintedstyle{tango}
\item Пространства имён могут быть вложенными:
\usemintedstyle{lightbulb}
\begin{tcblisting}{listing engine=minted, minted language=cpp, minted options={fontsize=\scriptsize}, listing only, colback=codebg, colframe=codebg, rounded corners}
namespace boost {
  namespace multiprecision {
    class gmp_int;
  }
}

boost::multiprecision::gmp_int x(42);
\end{tcblisting}
\pause
\usemintedstyle{tango}
\item Для пространств имён можно создавать новые имена для удобства (это гораздо лучше, чем \mintinline{cpp}|using namespace|):
\usemintedstyle{lightbulb}
\begin{tcblisting}{listing engine=minted, minted language=cpp, minted options={fontsize=\scriptsize}, listing only, colback=codebg, colframe=codebg, rounded corners}
namespace bmp = boost::multiprecision;

bmp::gmp_int x(42);
\end{tcblisting}
\end{itemize}
\end{frame}

\begin{frame}[fragile]
\frametitle{Глобальное пространство имён}
\begin{itemize}
\usemintedstyle{tango}
\item Технически, всё, что не лежит в каком-то пространстве имён, лежит в неявном \textit{глобальном пространстве имён}
\pause
\item Обратиться к объектам из глобального пространства имён можно через оператор разрешения области видимости \mintinline{cpp}|::|
\usemintedstyle{lightbulb}
\begin{tcblisting}{listing engine=minted, minted language=cpp, minted options={fontsize=\scriptsize}, listing only, colback=codebg, colframe=codebg, rounded corners}
namespace complex {
  class Complex { ... };

  Complex exp(Complex c) {
    float x = c.real;
    float y = c.imag;

    // ::exp(x) явно обращается к глобальному пространству имён
    return ::exp(x) * ...;
  }
}
\end{tcblisting}
\end{itemize}
\end{frame}

\begin{frame}[fragile]
\frametitle{Анонимное пространство имён}
\begin{itemize}
\usemintedstyle{tango}
\item У пространства имён можно опустить имя: тогда оно будет \textit{анонимным}
\pause
\item Объекты из него автоматически попадают в родительское пространство имён, но используют \textit{internal linkage} -- не видны из других cpp-файлов при линковке
\pause
\item То же самое, что и \mintinline{cpp}|static|, но удобнее, когда есть много объектов, и работает для классов:
\usemintedstyle{lightbulb}
\begin{tcblisting}{listing engine=minted, minted language=cpp, minted options={fontsize=\scriptsize}, listing only, colback=codebg, colframe=codebg, rounded corners}
static void helper();
static int global = 0;

namespace {
  void helper2();
  int global2 = 0;

  class Internal { ... };
}
\end{tcblisting}
\pause
\usemintedstyle{tango}
\item Совет: все детали реализации внутри cpp-файла прячьте в анонимное пространство имён
\end{itemize}
\end{frame}

\begin{frame}[fragile]
\frametitle{Inline пространство имён}
\begin{itemize}
\usemintedstyle{tango}
\item Пространство имён можно пометить как \mintinline{cpp}|inline|
\pause
\item Объекты из него автоматически попадают в родительское пространство имён, но всё равно используют имя пространства имён при name mangling'е
\pause
\item Иногда используется для версионирования библиотек (напр. libstdc++ -- стандартная библиотека C++ компилятора GCC)
\usemintedstyle{lightbulb}
\begin{tcblisting}{listing engine=minted, minted language=cpp, minted options={fontsize=\scriptsize}, listing only, colback=codebg, colframe=codebg, rounded corners}
namespace mylib {
  // Старый скомпилированный код использует v1
  namespace v1 {
    ...
  }

  // Новый код по умолчанию использует v2
  inline namespace v2 {
    ...
  }
}
\end{tcblisting}
\end{itemize}
\end{frame}

\begin{frame}[fragile]
\frametitle{Синонимы типов}
\begin{itemize}
\usemintedstyle{tango}
\item C-style синонимы типов \mintinline{cpp}|typedef int ObjectID;| часто тяжело воспринимаются
\pause
\item В C++ есть альтернативный синтаксис: \mintinline{cpp}|using ObjectID = int;|
\pause
\item Синонимы типов (как через \mintinline{cpp}|typedef|, так и через \mintinline{cpp}|using|) могут быть внутри классов или пространств имён:
\usemintedstyle{lightbulb}
\begin{tcblisting}{listing engine=minted, minted language=cpp, minted options={fontsize=\scriptsize}, listing only, colback=codebg, colframe=codebg, rounded corners}
namespace internal {
  using Data = int;
  struct Node {
    Data data;
    Node* next;
  };
}

class List {
  using Data = internal::Data;
  using Node = internal::Node;
  ...
};
\end{tcblisting}
\end{itemize}
\end{frame}

\begin{frame}[fragile]
\frametitle{Стандартная библиотека C++}
\begin{itemize}
\usemintedstyle{tango}
\item Алгоритмы: \mintinline{cpp}|<algorithm>|
\item Структуры данных: \mintinline{cpp}|<vector>, <list>, <forward_list>,| \mintinline{cpp}|<deque>, <set>, <map>, <unordered_set>, <unordered_map>|
\item Строки, поиск и форматирование: \mintinline{cpp}|<string>, <regex>,| \mintinline{cpp}|<format>, <print>|
\item Математика: \mintinline{cpp}|<cmath>, <complex>, <numeric>|
\item Генерация случайных чисел: \mintinline{cpp}|<random>|
\item Ввод/вывод: \mintinline{cpp}|<iostream>, <fstream>, <sstream>|
\item Работа с памятью, умные указатели: \mintinline{cpp}|<new>, <memory>,| \mintinline{cpp}|<memory_resource>|
\item Работа с временем: \mintinline{cpp}|<chrono>|
\item Многопоточное программирование: \mintinline{cpp}|<thread>, <mutex>,| \mintinline{cpp}|<future>, <condition_variable>|
\item И много другого
\end{itemize}
\end{frame}

\begin{frame}[fragile]
\frametitle{Стандартная библиотека C++}
\begin{itemize}
\usemintedstyle{tango}
\item Поддерживается любым компилятором C++
\pause
\item Все идентификаторы находятся в пространстве имён \mintinline{cpp}|std|, напр. \mintinline{cpp}|std::string|, \mintinline{cpp}|std::regex|, \mintinline{cpp}|std::format|
\pause
\item Все имена используют \mintinline{cpp}|snake_case|
\pause
\item Заголовочные файлы не имеют расширения: \mintinline{cpp}|<algorithm>|, \mintinline{cpp}|<future>|
\end{itemize}
\end{frame}

\begin{frame}[fragile]
\frametitle{Стандартная библиотека C++}
\begin{itemize}
\usemintedstyle{tango}
\item Заголовочные файлы стандартной библиотеки C тоже доступны в C++: \mintinline{cpp}|<string.h>|, \mintinline{cpp}|<math.h>| и т.д.
\pause
\item В стандартной библиотеке C++ есть файлы, оборачивающие C-файлы в \mintinline{cpp}|namespace std| -- рекомендуется использовать именно их
\pause
\item \mintinline{cpp}|<XXX.h>| превращается в \mintinline{cpp}|<cXXX>|: \mintinline{cpp}|<cstring>|, \mintinline{cpp}|<cmath>|, и т.д.
\pause
\item Например, функция \mintinline{cpp}|sin(float)| из \mintinline{cpp}|<math.h>| превращается в функцию \mintinline{cpp}|std::sin(float)| из \mintinline{cpp}|<cmath>|
\pause
\item Формально, \mintinline{cpp}|<cmath>| не обязан содержать функцию \mintinline{cpp}|sin| (без пространства имён \mintinline{cpp}|std|), но на практике часто содержит (но это не гарантируется!)
\pause
\item То же самое касается типов: \mintinline{cpp}|<cstdint>| содержит тип \mintinline{cpp}|std::uint32_t| и не обязан содержать тип \mintinline{cpp}|uint32_t|, но на практике часто содержит, и многие не пишут \mintinline{cpp}|std::| перед такими типами
\end{itemize}
\end{frame}

\begin{frame}[fragile]
\frametitle{Пример: C-style сортировка}
\begin{itemize}
\usemintedstyle{lightbulb}
\begin{tcblisting}{listing engine=minted, minted language=cpp, minted options={fontsize=\scriptsize}, listing only, colback=codebg, colframe=codebg, rounded corners}
struct point {
  int x, y;
};

void qsort(void* array, size_t count, size_t size,
  int (*comp)(const void*, const void*));

int cmp_point(const void* v1, const void* v2) {
  const struct point_s *p1 = (const struct point_s*)v1;
  const struct point_s *p2 = (const struct point_s*)v2;
  ...
}

struct point points[10];
qsort(points, 10, sizeof(points[0]), cmp_point);
\end{tcblisting}
\pause
\usemintedstyle{tango}
\item Можно использовать и для примитивных типов
\pause
\item Нет проверки типов со стороны компилятора (\mintinline{cpp}|void*|)
\pause
\item Функцию-компаратор \mintinline{cpp}|cmp_point| невозможно заинлайнить
\end{itemize}
\end{frame}

\begin{frame}[fragile]
\frametitle{Пример: ООП-style сортировка}
\begin{itemize}
\usemintedstyle{lightbulb}
\begin{tcblisting}{listing engine=minted, minted language=cpp, minted options={fontsize=\scriptsize}, listing only, colback=codebg, colframe=codebg, rounded corners}
class Comparable {
public:
  virtual int compare(const Comparable* other) const = 0;
};

void sort(Comparable** array, size_t count);

class Point : public Comparable {
private:
  int x, y;
public:
  int compare(const Comparable* other) const override {
    const Point* otherPoint = (const Point*)other;
    ...
  }
};

Point** points = new Point*[10];
...
nsort(points, 10);
\end{tcblisting}
\end{itemize}
\end{frame}

\begin{frame}[fragile]
\frametitle{Пример: ООП-style сортировка}
\begin{itemize}
\usemintedstyle{tango}
\item Проверка типов есть, но частичная -- всё ещё приходится делать приведение типов \mintinline{cpp}|Comparable* -> Point*|
\pause
\item Метод-компаратор \mintinline{cpp}|compare| всё ещё невозможно заинлайнить
\pause
\item Для примитивных типов приходится писать обёртки
\usemintedstyle{lightbulb}
\begin{tcblisting}{listing engine=minted, minted language=cpp, minted options={fontsize=\scriptsize}, listing only, colback=codebg, colframe=codebg, rounded corners}
class Integer : public Comparable {
private:
  int value;
public:
  int compare(const Comparable* other) const override {
    ...
  }
};
\end{tcblisting}
\end{itemize}
\end{frame}

\begin{frame}[fragile]
\frametitle{Пример: ООП-style сортировка}
\begin{itemize}
\usemintedstyle{tango}
\item Приходится хранить объекты по указателю
\pause
\item Сработает ли такой код?
\usemintedstyle{lightbulb}
\begin{tcblisting}{listing engine=minted, minted language=cpp, minted options={fontsize=\scriptsize}, listing only, colback=codebg, colframe=codebg, rounded corners}
void sort(Comparable* array, size_t count);

Point points[10];
sort(points, 10);
\end{tcblisting}
\pause
\usemintedstyle{tango}
\item Нет: при итерации по массиву функции \mintinline{cpp}|sort| нужно знать настоящий размер объектов (\mintinline{cpp}|sizeof(Point) > sizeof(Comparable)|)
\end{itemize}
\end{frame}

\begin{frame}[fragile]
\frametitle{Пример: C-style дерево поиска}
\begin{itemize}
\usemintedstyle{lightbulb}
\begin{tcblisting}{listing engine=minted, minted language=cpp, minted options={fontsize=\scriptsize}, listing only, colback=codebg, colframe=codebg, rounded corners}
struct node {
  void* data;
  struct node *left, *right;
};

struct tree {
  struct node* root;
};

typedef int (*cmp_fun)(const void*, const void*);

void insert(struct tree* t, void* data, cmp_fun cmp) {
  struct node* n = malloc(sizeof(struct node));
  ...
}

struct list l; l.head = NULL;
int* data = malloc(sizeof(int));
*data = 42;
insert(&l, data, cmp_int);
\end{tcblisting}
\pause
\usemintedstyle{tango}
\item Нет проверки типов, функцию-компаратор не заинлайнить
\pause
\item Две аллокации памяти!
\end{itemize}
\end{frame}

\begin{frame}[fragile]
\frametitle{Пример: ООП-style дерево поиска}
\begin{itemize}
\usemintedstyle{lightbulb}
\begin{tcblisting}{listing engine=minted, minted language=cpp, minted options={fontsize=\scriptsize}, listing only, colback=codebg, colframe=codebg, rounded corners}
class Tree {
  struct Node {
    Comparable* data;
    Node *left, *right;
  };

  Node* root;

public:

  void insert(Comparable* data) {
    Node* node = new Node;
    ...
  }
};

Tree tree;
tree.insert(new Point(3, 4));
\end{tcblisting}
\pause
\usemintedstyle{tango}
\item Все те же проблемы, что и с сортировкой
\end{itemize}
\end{frame}

\begin{frame}[fragile]
\frametitle{Шаблоны}
\begin{itemize}
\usemintedstyle{tango}
\item Все проблемы в примерах исчезли бы, если бы мы умели \textit{параметризовать} фукнции/классы некоторыми \textit{типами}
\pause
\item Сейчас мы умеем так:
\usemintedstyle{lightbulb}
\begin{tcblisting}{listing engine=minted, minted language=cpp, minted options={fontsize=\scriptsize}, listing only, colback=codebg, colframe=codebg, rounded corners}
// Приходится писать новый класс для
// каждого хранимого типа данных
class Array {
  int* data;
  size_t size;
};

// Приходится писать новый класс для
// каждого хранимого типа данных
class scoped_ptr {
  int* data;
};
\end{tcblisting}
\end{itemize}
\end{frame}

\begin{frame}[fragile]
\frametitle{Шаблоны}
\begin{itemize}
\usemintedstyle{tango}
\item А хочется как-то так:
\usemintedstyle{lightbulb}
\begin{tcblisting}{listing engine=minted, minted language=cpp, minted options={fontsize=\scriptsize}, listing only, colback=codebg, colframe=codebg, rounded corners}
class Array {
  // T - как-то заданный, но произвольный тип
  T* data;
  size_t size;
};

class scoped_ptr {
  // T - как-то заданный, но произвольный тип
  T* data;
};
\end{tcblisting}
\end{itemize}
\end{frame}

\begin{frame}[fragile]
\frametitle{Шаблоны на C препроцессоре}
\begin{itemize}
\usemintedstyle{tango}
\item Решение 1: использовать препроцессор
\usemintedstyle{lightbulb}
\begin{tcblisting}{listing engine=minted, minted language=cpp, minted options={fontsize=\scriptsize}, listing only, colback=codebg, colframe=codebg, rounded corners}
#define Array(T) class Array_##T { \
private: \
  T* data; \
  size_t size; \
public: \
  T& operator[](size_t index) { \
    return array[index]; \
  } \
  ... \
};

Array(int) // генерируем код класса Array_int
Array(float) // генерируем код класса Array_float

Array_int array(10);
array[5] = 42;
\end{tcblisting}
\end{itemize}
\end{frame}

\begin{frame}[fragile]
\frametitle{Шаблоны на C препроцессоре}
\begin{itemize}
\usemintedstyle{tango}
\item Есть все проблемы от препроцессора: простая текстовая подстановка, неудобство отладки, и т.п.
\pause
\item Классы \mintinline{cpp}|Array_int| и \mintinline{cpp}|Array_float| никак не связаны между собой
\end{itemize}
\end{frame}

\begin{frame}[fragile]
\frametitle{Шаблоны в C++}
\begin{itemize}
\usemintedstyle{tango}
\item Решение 2: использовать C++ шаблоны
\pause
\usemintedstyle{lightbulb}
\begin{tcblisting}{listing engine=minted, minted language=cpp, minted options={fontsize=\scriptsize}, listing only, colback=codebg, colframe=codebg, rounded corners}
template <typename T>
class Array {
private:
  T* data;
  size_t size;
public:
  T& operator[](size_t index) {
    return array[index];
  }
  ...
};

Array<int> array(10);
array[5] = 42;
\end{tcblisting}
\end{itemize}
\end{frame}

\begin{frame}[fragile]
\frametitle{Шаблоны в C++}
\begin{itemize}
\usemintedstyle{tango}
\item Подстановку делает не препроцессор, а компилятор -- он знает про типы, про области видимости, и про другие конструкции языка
\pause
\item Работает посредством \textit{мономорфизации}: компилятор внутри фактически генерирует отдельные классы \mintinline{cpp}|Array<int>|, \mintinline{cpp}|Array<float>| и т.д.
\pause
\item Код шаблонного класса проще читать, и можно разбить на несколько файлов при желании (но всё равно должен быть в заголовочном файле)
\end{itemize}
\end{frame}

\begin{frame}[fragile]
\frametitle{Шаблоны в C++: примеры}
\begin{itemize}
\usemintedstyle{lightbulb}
\begin{multicols}{2}
\begin{tcblisting}{listing engine=minted, minted language=cpp, minted options={fontsize=\scriptsize}, listing only, colback=codebg, colframe=codebg, rounded corners}
template <typename T>
class scoped_ptr {
  T* data;
  ...
};
\end{tcblisting}
\begin{tcblisting}{listing engine=minted, minted language=cpp, minted options={fontsize=\scriptsize}, listing only, colback=codebg, colframe=codebg, rounded corners}
template <typename T>
class list {
  struct node {
    T data;
    node* next;
  };

  node* head;
  ...
};
\end{tcblisting}
\columnbreak
\begin{tcblisting}{listing engine=minted, minted language=cpp, minted options={fontsize=\scriptsize}, listing only, colback=codebg, colframe=codebg, rounded corners}
template <typename T>
class tree {
  struct node {
    T data;
    node *left, *right;
  };

  node* root;
  ...
};
\end{tcblisting}
\begin{tcblisting}{listing engine=minted, minted language=cpp, minted options={fontsize=\scriptsize}, listing only, colback=codebg, colframe=codebg, rounded corners}
template <typename T>
void sort(T* array,
  size_t count) {
  ...
}
\end{tcblisting}
\end{multicols}
\end{itemize}
\end{frame}

\end{document}